% Options for packages loaded elsewhere
\PassOptionsToPackage{unicode}{hyperref}
\PassOptionsToPackage{hyphens}{url}
%
\documentclass[
  10pt,
  ignorenonframetext,
  aspectratio=169,
]{beamer}
\usepackage{pgfpages}
\setbeamertemplate{caption}[numbered]
\setbeamertemplate{caption label separator}{: }
\setbeamercolor{caption name}{fg=normal text.fg}
\beamertemplatenavigationsymbolshorizontal
% Prevent slide breaks in the middle of a paragraph
\widowpenalties 1 10000
\raggedbottom
\setbeamertemplate{part page}{
  \centering
  \begin{beamercolorbox}[sep=16pt,center]{part title}
    \usebeamerfont{part title}\insertpart\par
  \end{beamercolorbox}
}
\setbeamertemplate{section page}{
  \centering
  \begin{beamercolorbox}[sep=12pt,center]{part title}
    \usebeamerfont{section title}\insertsection\par
  \end{beamercolorbox}
}
\setbeamertemplate{subsection page}{
  \centering
  \begin{beamercolorbox}[sep=8pt,center]{part title}
    \usebeamerfont{subsection title}\insertsubsection\par
  \end{beamercolorbox}
}
\AtBeginPart{
  \frame{\partpage}
}
\AtBeginSection{
  \ifbibliography
  \else
    \frame{\sectionpage}
  \fi
}
\AtBeginSubsection{
  \frame{\subsectionpage}
}

\usepackage{amsmath,amssymb}
\usepackage{iftex}
\ifPDFTeX
  \usepackage[T1]{fontenc}
  \usepackage[utf8]{inputenc}
  \usepackage{textcomp} % provide euro and other symbols
\else % if luatex or xetex
  \usepackage{unicode-math}
  \defaultfontfeatures{Scale=MatchLowercase}
  \defaultfontfeatures[\rmfamily]{Ligatures=TeX,Scale=1}
\fi
\usepackage{lmodern}
\usetheme[]{Madrid}
\usecolortheme{seahorse}
\usefonttheme{structurebold}
\ifPDFTeX\else  
    % xetex/luatex font selection
\fi
% Use upquote if available, for straight quotes in verbatim environments
\IfFileExists{upquote.sty}{\usepackage{upquote}}{}
\IfFileExists{microtype.sty}{% use microtype if available
  \usepackage[]{microtype}
  \UseMicrotypeSet[protrusion]{basicmath} % disable protrusion for tt fonts
}{}
\makeatletter
\@ifundefined{KOMAClassName}{% if non-KOMA class
  \IfFileExists{parskip.sty}{%
    \usepackage{parskip}
  }{% else
    \setlength{\parindent}{0pt}
    \setlength{\parskip}{6pt plus 2pt minus 1pt}}
}{% if KOMA class
  \KOMAoptions{parskip=half}}
\makeatother
\usepackage{xcolor}
\newif\ifbibliography
\setlength{\emergencystretch}{3em} % prevent overfull lines
\setcounter{secnumdepth}{-\maxdimen} % remove section numbering


\providecommand{\tightlist}{%
  \setlength{\itemsep}{0pt}\setlength{\parskip}{0pt}}\usepackage{longtable,booktabs,array}
\usepackage{calc} % for calculating minipage widths
\usepackage{caption}
% Make caption package work with longtable
\makeatletter
\def\fnum@table{\tablename~\thetable}
\makeatother
\usepackage{graphicx}
\makeatletter
\def\maxwidth{\ifdim\Gin@nat@width>\linewidth\linewidth\else\Gin@nat@width\fi}
\def\maxheight{\ifdim\Gin@nat@height>\textheight\textheight\else\Gin@nat@height\fi}
\makeatother
% Scale images if necessary, so that they will not overflow the page
% margins by default, and it is still possible to overwrite the defaults
% using explicit options in \includegraphics[width, height, ...]{}
\setkeys{Gin}{width=\maxwidth,height=\maxheight,keepaspectratio}
% Set default figure placement to htbp
\makeatletter
\def\fps@figure{htbp}
\makeatother

\usepackage{booktabs}
\usepackage{longtable}
\usepackage{array}
\usepackage{multirow}
\usepackage{wrapfig}
\usepackage{float}
\usepackage{colortbl}
\usepackage{pdflscape}
\usepackage{tabu}
\usepackage{threeparttable}
\usepackage{threeparttablex}
\usepackage[normalem]{ulem}
\usepackage{makecell}
\usepackage{xcolor}
\usepackage{booktabs}
\usepackage{longtable}
\usepackage{array}
\usepackage{multirow}
\usepackage{xcolor}
\definecolor{churnYes}{HTML}{E24B4A}
\definecolor{churnNo}{HTML}{4999DA}
\definecolor{accent}{HTML}{185FA5}
\setbeamertemplate{caption}[numbered]
\setbeamertemplate{itemize items}[circle]
\setbeamertemplate{enumerate items}[default]
\makeatletter
\@ifpackageloaded{caption}{}{\usepackage{caption}}
\AtBeginDocument{%
\ifdefined\contentsname
  \renewcommand*\contentsname{Table of contents}
\else
  \newcommand\contentsname{Table of contents}
\fi
\ifdefined\listfigurename
  \renewcommand*\listfigurename{List of Figures}
\else
  \newcommand\listfigurename{List of Figures}
\fi
\ifdefined\listtablename
  \renewcommand*\listtablename{List of Tables}
\else
  \newcommand\listtablename{List of Tables}
\fi
\ifdefined\figurename
  \renewcommand*\figurename{Figure}
\else
  \newcommand\figurename{Figure}
\fi
\ifdefined\tablename
  \renewcommand*\tablename{Table}
\else
  \newcommand\tablename{Table}
\fi
}
\@ifpackageloaded{float}{}{\usepackage{float}}
\floatstyle{ruled}
\@ifundefined{c@chapter}{\newfloat{codelisting}{h}{lop}}{\newfloat{codelisting}{h}{lop}[chapter]}
\floatname{codelisting}{Listing}
\newcommand*\listoflistings{\listof{codelisting}{List of Listings}}
\makeatother
\makeatletter
\makeatother
\makeatletter
\@ifpackageloaded{caption}{}{\usepackage{caption}}
\@ifpackageloaded{subcaption}{}{\usepackage{subcaption}}
\makeatother
\ifLuaTeX
  \usepackage{selnolig}  % disable illegal ligatures
\fi
\usepackage{bookmark}

\IfFileExists{xurl.sty}{\usepackage{xurl}}{} % add URL line breaks if available
\urlstyle{same} % disable monospaced font for URLs
\hypersetup{
  pdftitle={Prédiction du Churn Client},
  pdfauthor={Master MECEN --- Semestre 2},
  hidelinks,
  pdfcreator={LaTeX via pandoc}}

\title{Prédiction du Churn Client}
\subtitle{Classification supervisée avec tidymodels}
\author{Master MECEN --- Semestre 2}
\date{29 March 2026}

\begin{document}
\frame{\titlepage}

\renewcommand*\contentsname{Table of contents}
\begin{frame}[allowframebreaks]
  \frametitle{Table of contents}
  \tableofcontents[hideallsubsections]
\end{frame}
\section{Introduction}\label{introduction}

\begin{frame}{Contexte et objectif}
\phantomsection\label{contexte-et-objectif}
\begin{block}{Problématique métier}
\phantomsection\label{probluxe9matique-muxe9tier}
Le \textbf{churn} (désabonnement) représente un enjeu majeur pour les
opérateurs télécoms :

\begin{itemize}
\tightlist
\item
  Coût d'acquisition d'un nouveau client : \textbf{5 à 25×} le coût de
  rétention
\item
  Identifier les clients à risque permet des \textbf{actions préventives
  ciblées}
\item
  Optimisation des ressources marketing et fidélisation
\end{itemize}
\end{block}

\begin{block}{Objectif de l'étude}
\phantomsection\label{objectif-de-luxe9tude}
\begin{quote}
Développer un modèle prédictif performant pour identifier les clients
susceptibles de résilier leur abonnement.
\end{quote}

\textbf{Source} : Iranian Churn Dataset (UCI Machine Learning
Repository)
\end{block}
\end{frame}

\begin{frame}[fragile]{Méthodologie adoptée}
\phantomsection\label{muxe9thodologie-adoptuxe9e}
\begin{block}{Pipeline de modélisation}
\phantomsection\label{pipeline-de-moduxe9lisation}
\begin{center}
\textbf{EDA} $\rightarrow$ \textbf{Preprocessing} $\rightarrow$ \textbf{Benchmark} $\rightarrow$ \textbf{Tuning} $\rightarrow$ \textbf{Évaluation}
\end{center}
\end{block}

\begin{block}{Framework technique}
\phantomsection\label{framework-technique}
\begin{itemize}
\tightlist
\item
  \textbf{Langage} : R avec l'écosystème \texttt{tidymodels}
\item
  \textbf{Packages} : \texttt{recipes}, \texttt{parsnip},
  \texttt{workflows}, \texttt{tune}, \texttt{yardstick}, \texttt{themis}
\item
  \textbf{Validation} : Cross-validation stratifiée (10 folds)
\item
  \textbf{Métriques} : ROC AUC (principale), F1-score, Precision, Recall
\end{itemize}
\end{block}
\end{frame}

\section{Présentation des données}\label{pruxe9sentation-des-donnuxe9es}

\begin{frame}{Vue d'ensemble du dataset}
\phantomsection\label{vue-densemble-du-dataset}
\begingroup\fontsize{9}{11}\selectfont

\begin{longtable*}[t]{lc}
\toprule
Caractéristique & Valeur\\
\midrule
Observations & 3 150\\
Variables & 14\\
Variable cible & Churn (binaire)\\
Numériques & 8\\
Catégorielles & 6\\
\addlinespace
Valeurs manquantes & 0\\
\bottomrule
\end{longtable*}
\endgroup{}

\begin{block}{Caractéristiques}
\phantomsection\label{caractuxe9ristiques}
Données issues d'un opérateur télécom iranien :

\begin{itemize}
\tightlist
\item
  \textbf{Comportement d'usage} : appels, SMS, durée d'utilisation
\item
  \textbf{Profil client} : âge, ancienneté, valeur client
\item
  \textbf{Relation commerciale} : plaintes, forfait, statut
\end{itemize}
\end{block}
\end{frame}

\begin{frame}{Dictionnaire des variables}
\phantomsection\label{dictionnaire-des-variables}
\begingroup\fontsize{7}{9}\selectfont

\begin{longtable*}[t]{llc}
\toprule
 & Type & Description\\
\midrule
\cellcolor{gray!10}{Call  Failure} & \cellcolor{gray!10}{Entier} & \cellcolor{gray!10}{Nombre d'appels échoués}\\
Complains & Facteur & Indicateur de plainte\\
\cellcolor{gray!10}{Subscription  Length} & \cellcolor{gray!10}{Entier} & \cellcolor{gray!10}{Durée abonnement (mois)}\\
Charge  Amount & Ordonné & Montant facturé (0-10)\\
\cellcolor{gray!10}{Seconds of Use} & \cellcolor{gray!10}{Entier} & \cellcolor{gray!10}{Secondes d'appels}\\
\addlinespace
Frequency of use & Entier & Nombre total d'appels\\
\cellcolor{gray!10}{Frequency of SMS} & \cellcolor{gray!10}{Entier} & \cellcolor{gray!10}{Nombre de SMS}\\
Distinct Called Numbers & Entier & Numéros uniques appelés\\
\cellcolor{gray!10}{Age Group} & \cellcolor{gray!10}{Ordonné} & \cellcolor{gray!10}{Tranche d'âge (1-5)}\\
Tariff Plan & Facteur & Type de forfait\\
\addlinespace
\cellcolor{gray!10}{Status} & \cellcolor{gray!10}{Facteur} & \cellcolor{gray!10}{Statut du compte}\\
Age & Entier & Âge du client\\
\cellcolor{gray!10}{Customer Value} & \cellcolor{gray!10}{Numérique} & \cellcolor{gray!10}{Valeur client estimée}\\
\textcolor[HTML]{E24B4A}{\textbf{Churn}} & \textcolor[HTML]{E24B4A}{\textbf{Facteur}} & \textcolor[HTML]{E24B4A}{\textbf{Cible : churn}}\\
\bottomrule
\end{longtable*}
\endgroup{}
\end{frame}

\section{Analyse exploratoire}\label{analyse-exploratoire}

\begin{frame}{Déséquilibre de la variable cible}
\phantomsection\label{duxe9suxe9quilibre-de-la-variable-cible}
\begin{center}
\includegraphics{Presentation_files/figure-beamer/desequilibre-1.pdf}
\end{center}

\begin{block}{Conséquence méthodologique}
\phantomsection\label{consuxe9quence-muxe9thodologique}
\begin{itemize}
\tightlist
\item
  \textbf{Accuracy trompeuse} : modèle naïf → 84\%
\item
  \textbf{Solution} : SMOTE pour rééquilibrer
\item
  \textbf{Métrique principale} : ROC AUC
\end{itemize}
\end{block}
\end{frame}

\begin{frame}{Distribution des variables numériques}
\phantomsection\label{distribution-des-variables-numuxe9riques}
\end{frame}

\section{\texorpdfstring{\texttt{\{r\ summary-num\}\ \#\ data\ \textbar{}\textgreater{}\ \#\ \ \ select(where(is.numeric))\ \textbar{}\textgreater{}\ \#\ \ \ pivot\_longer(everything(),\ names\_to\ =\ "Variable")\ \textbar{}\textgreater{}\ \#\ \ \ group\_by(Variable)\ \textbar{}\textgreater{}\ \#\ \ \ summarise(\ \#\ \ \ \ \ Min\ \ \ \ \ =\ min(value,\ na.rm\ =\ TRUE),\ \#\ \ \ \ \ Médiane\ =\ median(value,\ na.rm\ =\ TRUE),\ \#\ \ \ \ \ Moyenne\ =\ round(mean(value,\ na.rm\ =\ TRUE),\ 1),\ \#\ \ \ \ \ Max\ \ \ \ \ =\ max(value,\ na.rm\ =\ TRUE),\ \#\ \ \ \ \ \textasciigrave{}E.T.\textasciigrave{}\ \ =\ round(sd(value,\ na.rm\ =\ TRUE),\ 1)\ \#\ \ \ )\ \textbar{}\textgreater{}\ \#\ \ \ kable(booktabs\ =\ TRUE,\ align\ =\ c("l",\ rep("r",\ 5)))\ \textbar{}\textgreater{}\ \#\ \ \ kable\_styling(font\_size\ =\ 8,\ latex\_options\ =\ c("striped",\ "hold\_position"))}}{\{r summary-num\} \# data \textbar\textgreater{} \#   select(where(is.numeric)) \textbar\textgreater{} \#   pivot\_longer(everything(), names\_to = "Variable") \textbar\textgreater{} \#   group\_by(Variable) \textbar\textgreater{} \#   summarise( \#     Min     = min(value, na.rm = TRUE), \#     Médiane = median(value, na.rm = TRUE), \#     Moyenne = round(mean(value, na.rm = TRUE), 1), \#     Max     = max(value, na.rm = TRUE), \#     `E.T.`  = round(sd(value, na.rm = TRUE), 1) \#   ) \textbar\textgreater{} \#   kable(booktabs = TRUE, align = c("l", rep("r", 5))) \textbar\textgreater{} \#   kable\_styling(font\_size = 8, latex\_options = c("striped", "hold\_position"))}}\label{r-summary-num-data-selectwhereis.numeric-pivot_longereverything-names_to-variable-group_byvariable-summarise-min-minvalue-na.rm-true-muxe9diane-medianvalue-na.rm-true-moyenne-roundmeanvalue-na.rm-true-1-max-maxvalue-na.rm-true-e.t.-roundsdvalue-na.rm-true-1-kablebooktabs-true-align-cl-repr-5-kable_stylingfont_size-8-latex_options-cstriped-hold_position}

\begin{frame}{\texttt{\{r\ summary-num\}\ \#\ data\ \textbar{}\textgreater{}\ \#\ \ \ select(where(is.numeric))\ \textbar{}\textgreater{}\ \#\ \ \ pivot\_longer(everything(),\ names\_to\ =\ "Variable")\ \textbar{}\textgreater{}\ \#\ \ \ group\_by(Variable)\ \textbar{}\textgreater{}\ \#\ \ \ summarise(\ \#\ \ \ \ \ Min\ \ \ \ \ =\ min(value,\ na.rm\ =\ TRUE),\ \#\ \ \ \ \ Médiane\ =\ median(value,\ na.rm\ =\ TRUE),\ \#\ \ \ \ \ Moyenne\ =\ round(mean(value,\ na.rm\ =\ TRUE),\ 1),\ \#\ \ \ \ \ Max\ \ \ \ \ =\ max(value,\ na.rm\ =\ TRUE),\ \#\ \ \ \ \ \textasciigrave{}E.T.\textasciigrave{}\ \ =\ round(sd(value,\ na.rm\ =\ TRUE),\ 1)\ \#\ \ \ )\ \textbar{}\textgreater{}\ \#\ \ \ kable(booktabs\ =\ TRUE,\ align\ =\ c("l",\ rep("r",\ 5)))\ \textbar{}\textgreater{}\ \#\ \ \ kable\_styling(font\_size\ =\ 8,\ latex\_options\ =\ c("striped",\ "hold\_position"))}}
\end{frame}

\begin{frame}{Variables catégorielles --- Taux de churn}
\phantomsection\label{variables-catuxe9gorielles-taux-de-churn}
\begin{center}
\includegraphics{Presentation_files/figure-beamer/cat-churn-1.pdf}
\end{center}
\end{frame}

\begin{frame}{Synthèse EDA --- Variables clés}
\phantomsection\label{synthuxe8se-eda-variables-cluxe9s}
\begin{block}{Variables à fort pouvoir discriminant}
\phantomsection\label{variables-uxe0-fort-pouvoir-discriminant}
\begin{longtable}[]{@{}lll@{}}
\toprule\noalign{}
Variable & Observation & Implication \\
\midrule\noalign{}
\endhead
\textbf{Complains} & Beaucoup de plainte → churn & Prédicteur majeur \\
\textbf{Status} & Comptes inactifs & Signal fort \\
\textbf{Frequency of use} & Faible usage → churn & Engagement \\
\textbf{Customer Value} & Valeurs extrêmes & Segmentation \\
\bottomrule\noalign{}
\end{longtable}
\end{block}

\begin{block}{Points de vigilance}
\phantomsection\label{points-de-vigilance}
\begin{itemize}
\tightlist
\item
  \textbf{Multicolinéarité} : Frequency of use ↔ Seconds of Use (r ≈
  0.95)
\item
  \textbf{Variables ordinales} : Age Group, Charge Amount
\end{itemize}
\end{block}
\end{frame}

\section{Preprocessing}\label{preprocessing}

\begin{frame}[fragile]{Stratégie d'échantillonnage}
\phantomsection\label{stratuxe9gie-duxe9chantillonnage}
\begin{block}{Split train / test = 80\%}
\phantomsection\label{split-train-test-80}
\begingroup\fontsize{9}{11}\selectfont

\begin{longtable*}[t]{lrc}
\toprule
Ensemble & Observations & Ratio Churn = Yes\\
\midrule
Train & 2520 & 15.7\%\\
Test & 630 & 15.7\%\\
\bottomrule
\end{longtable*}
\endgroup{}
\end{block}

\begin{block}{Validation croisée}
\phantomsection\label{validation-croisuxe9e}
\begin{itemize}
\tightlist
\item
  \textbf{10 folds stratifiés} sur le train set
\item
  Stratification sur \texttt{Churn} → \textasciitilde16\% de ``Yes'' par
  fold
\item
  Estimations robustes et comparables
\end{itemize}
\end{block}
\end{frame}

\begin{frame}[fragile]{Recettes différenciées par famille}
\phantomsection\label{recettes-diffuxe9renciuxe9es-par-famille}
\begin{block}{Principe (cf.~tmwr.org/pre-proc-table)}
\phantomsection\label{principe-cf.-tmwr.orgpre-proc-table}
\begin{longtable}[]{@{}lll@{}}
\toprule\noalign{}
Recette & Modèles & Spécificités \\
\midrule\noalign{}
\endhead
\texttt{recipe\_tree} & RF, Bagging, DT & Pas de dummies ni
normalisation \\
\texttt{recipe\_xgb} & XGBoost & Dummies one-hot \\
\texttt{recipe\_distance} & KNN, SVM, Logit & Normalisation + dummies \\
\bottomrule\noalign{}
\end{longtable}
\end{block}

\begin{block}{Gestion du déséquilibre}
\phantomsection\label{gestion-du-duxe9suxe9quilibre}
\begin{itemize}
\tightlist
\item
  \textbf{SMOTENC} pour arbres (préserve les facteurs)
\item
  \textbf{SMOTE} après dummification pour les autres
\item
  Application en \textbf{dernière étape} uniquement
\end{itemize}
\end{block}
\end{frame}

\section{Benchmark initial}\label{benchmark-initial}

\begin{frame}{Modèles en compétition}
\phantomsection\label{moduxe8les-en-compuxe9tition}
\begin{block}{Arbres et ensemble}
\phantomsection\label{arbres-et-ensemble}
\begin{longtable}[]{@{}ll@{}}
\toprule\noalign{}
Modèle & Caractéristique \\
\midrule\noalign{}
\endhead
Decision Tree & Interprétable \\
Bagging & mtry = p \\
Random Forest & mtry = √p, 500 arbres \\
XGBoost & Boosting, early stopping \\
\bottomrule\noalign{}
\end{longtable}
\end{block}

\begin{block}{Distance}
\phantomsection\label{distance}
\begin{longtable}[]{@{}ll@{}}
\toprule\noalign{}
Modèle & Caractéristique \\
\midrule\noalign{}
\endhead
KNN & k = 10 \\
SVM linéaire & Frontière linéaire \\
SVM RBF & Kernel radial \\
\bottomrule\noalign{}
\end{longtable}
\end{block}
\end{frame}

\begin{frame}{Résultats --- ROC AUC}
\phantomsection\label{ruxe9sultats-roc-auc}
\begin{center}
\includegraphics[width=0.8\textwidth,height=\textheight]{Presentation_files/figure-beamer/benchmark-roc-1.pdf}
\end{center}
\end{frame}

\begin{frame}{Tableau comparatif}
\phantomsection\label{tableau-comparatif}
\begingroup\fontsize{8}{10}\selectfont

\begin{longtable*}[t]{lcccc}
\toprule
Modèle & F1 & Prec. & Recall & ROC AUC\\
\midrule
\cellcolor{gray!10}{xgb\_xgb} & \cellcolor{gray!10}{97.6\%} & \cellcolor{gray!10}{97.5\%} & \cellcolor{gray!10}{97.6\%} & \cellcolor{gray!10}{98.5\%}\\
tree\_rf & 97.5\% & 98.5\% & 96.5\% & 98.4\%\\
\cellcolor{gray!10}{tree\_bag} & \cellcolor{gray!10}{97.2\%} & \cellcolor{gray!10}{97.7\%} & \cellcolor{gray!10}{96.7\%} & \cellcolor{gray!10}{97.5\%}\\
dist\_knn & 96.3\% & 97.8\% & 94.9\% & 96.7\%\\
\cellcolor{gray!10}{dist\_svm\_rad} & \cellcolor{gray!10}{91.9\%} & \cellcolor{gray!10}{98.8\%} & \cellcolor{gray!10}{85.8\%} & \cellcolor{gray!10}{96.3\%}\\
\addlinespace
dist\_logit & 90\% & 97.1\% & 83.9\% & 93.2\%\\
\cellcolor{gray!10}{dist\_svm\_lin} & \cellcolor{gray!10}{90.5\%} & \cellcolor{gray!10}{96.9\%} & \cellcolor{gray!10}{84.8\%} & \cellcolor{gray!10}{92.9\%}\\
tree\_dt & 91.7\% & 98.1\% & 86.2\% & 91.7\%\\
\bottomrule
\end{longtable*}
\endgroup{}
\end{frame}

\begin{frame}{Courbes ROC}
\phantomsection\label{courbes-roc}
\begin{center}
\includegraphics[width=0.8\textwidth,height=\textheight]{Presentation_files/figure-beamer/roc-curves-1.pdf}
\end{center}
\end{frame}

\section{Focus XGBoost --- Tuning}\label{focus-xgboost-tuning}

\begin{frame}[fragile]{Stratégie en 2 phases}
\phantomsection\label{stratuxe9gie-en-2-phases}
\begin{block}{Phase 1 : Exploration large}
\phantomsection\label{phase-1-exploration-large}
\begin{itemize}
\tightlist
\item
  \textbf{Méthode} : Latin Hypercube (50 combinaisons)
\item
  \textbf{Algorithme} : \texttt{tune\_race\_anova()}
\end{itemize}
\end{block}

\begin{block}{Phase 2 : Zoom local}
\phantomsection\label{phase-2-zoom-local}
\begin{itemize}
\tightlist
\item
  \textbf{Méthode} : Grille resserrée (40 combinaisons)
\item
  \textbf{Validation} : CV 10-fold stratifié
\end{itemize}
\end{block}

\begin{block}{Hyperparamètres tunés}
\phantomsection\label{hyperparamuxe8tres-tunuxe9s}
\begin{longtable}[]{@{}ll@{}}
\toprule\noalign{}
Paramètre & Range \\
\midrule\noalign{}
\endhead
\texttt{trees} & 50 -- 900 \\
\texttt{tree\_depth} & 1 -- 6 \\
\texttt{learn\_rate} & 0.001 -- 0.1 \\
\texttt{min\_n} & 2 -- 20 \\
\bottomrule\noalign{}
\end{longtable}
\end{block}
\end{frame}

\begin{frame}{Visualisation du tuning}
\phantomsection\label{visualisation-du-tuning}
\begin{center}
\includegraphics[width=0.8\textwidth,height=\textheight]{Presentation_files/figure-beamer/xgb-tuning-1.pdf}
\end{center}
\end{frame}

\begin{frame}{Meilleurs hyperparamètres en phase 2}
\phantomsection\label{meilleurs-hyperparamuxe8tres-en-phase-2}
\begingroup\fontsize{9}{11}\selectfont

\begin{longtable*}[t]{lr}
\toprule
Paramètre & Valeur\\
\midrule
mtry & 7.0000\\
trees & 549.0000\\
min\_n & 5.0000\\
tree\_depth & 3.0000\\
learn\_rate & 0.0616\\
\addlinespace
loss\_reduction & 0.0173\\
sample\_size & 0.9192\\
\bottomrule
\end{longtable*}
\endgroup{}
\end{frame}

\begin{frame}{Importance des variables}
\phantomsection\label{importance-des-variables}
\begin{center}
\includegraphics[width=0.8\textwidth,height=\textheight]{Presentation_files/figure-beamer/xgb-importance-1.pdf}
\end{center}
\end{frame}

\section{Focus Random Forest}\label{focus-random-forest}

\begin{frame}[fragile]{Configuration du tuning}
\phantomsection\label{configuration-du-tuning}
\begin{block}{Grille testée}
\phantomsection\label{grille-testuxe9e}
\begin{longtable}[]{@{}ll@{}}
\toprule\noalign{}
Paramètre & Niveaux \\
\midrule\noalign{}
\endhead
\texttt{mtry} & 2, 4, 6, 9, 13 \\
\texttt{min\_n} & 2, 5, 11, 21, 40 \\
\bottomrule\noalign{}
\end{longtable}

\textbf{25 combinaisons} en CV 10-fold
\end{block}
\end{frame}

\begin{frame}[fragile]{Comparaison initial vs tuné}
\phantomsection\label{comparaison-initial-vs-tunuxe9}
\begingroup\fontsize{9}{11}\selectfont

\begin{longtable*}[t]{lcc>{}c}
\toprule
Métrique & Initial & Tuné & Gain\\
\midrule
ROC AUC & 97.5\% & 98\% & \textcolor[HTML]{1D9E75}{\textbf{+0.5\%}}\\
F1-Score & 98.5\% & 99.9\% & \textcolor[HTML]{1D9E75}{\textbf{+1.4\%}}\\
Precision & 96.5\% & 97.8\% & \textcolor[HTML]{1D9E75}{\textbf{+1.3\%}}\\
Recall & 98.4\% & 99.9\% & \textcolor[HTML]{1D9E75}{\textbf{+1.5\%}}\\
\bottomrule
\end{longtable*}
\endgroup{}

\begin{block}{Observations}
\phantomsection\label{observations}
\begin{itemize}
\tightlist
\item
  Gains modestes mais significatifs
\item
  Amélioration du \textbf{recall} (+1.5 pts)
\item
  \texttt{min\_n} élevé évite le surapprentissage
\end{itemize}
\end{block}
\end{frame}

\section{Modèle final}\label{moduxe8le-final}

\begin{frame}{XGBoost --- Évaluation test set}
\phantomsection\label{xgboost-uxe9valuation-test-set}
\begin{block}{Performances finales}
\phantomsection\label{performances-finales}
\begingroup\fontsize{9}{11}\selectfont

\begin{longtable*}[t]{lccc}
\toprule
Métrique & CV & Test & Delta\\
\midrule
f\_meas & 97.4\% & 97.4\% & \textasciitilde{}-1\%\\
precision & 98.3\% & 98.6\% & \textasciitilde{}-1\%\\
recall & 96.5\% & 96.2\% & \textasciitilde{}-1\%\\
roc\_auc & 98.7\% & 98.9\% & \textasciitilde{}-1\%\\
\bottomrule
\end{longtable*}
\endgroup{}
\end{block}
\end{frame}

\begin{frame}{Matrice de confusion}
\phantomsection\label{matrice-de-confusion}
\begin{center}
\includegraphics{Presentation_files/figure-beamer/conf-matrix-1.pdf}
\end{center}
\end{frame}

\begin{frame}{Distribution des probabilités}
\phantomsection\label{distribution-des-probabilituxe9s}
\begin{center}
\includegraphics{Presentation_files/figure-beamer/prob-dist-1.pdf}
\end{center}
\end{frame}

\begin{frame}{Limites et améliorations}
\phantomsection\label{limites-et-amuxe9liorations}
\begin{block}{Limites}
\phantomsection\label{limites}
\begin{longtable}[]{@{}ll@{}}
\toprule\noalign{}
Aspect & Limite \\
\midrule\noalign{}
\endhead
\textbf{Données} & Dataset monolingue \\
\textbf{Temporalité} & Pas de dimension temporelle \\
\textbf{Interprétabilité} & Boîte noire \\
\textbf{Coûts} & FN ≠ FP non intégrés \\
\bottomrule\noalign{}
\end{longtable}
\end{block}

\begin{block}{Pistes d'amélioration}
\phantomsection\label{pistes-damuxe9lioration}
\begin{enumerate}
\tightlist
\item
  Feature engineering temporel
\item
  Calibration (Platt scaling)
\item
  Coûts asymétriques
\item
  SHAP values
\item
  Stacking XGBoost + RF
\end{enumerate}
\end{block}
\end{frame}

\section{Conclusion}\label{conclusion}

\begin{frame}{Synthèse}
\phantomsection\label{synthuxe8se}
\begin{block}{Modèle recommandé : XGBoost}
\phantomsection\label{moduxe8le-recommanduxe9-xgboost}
\begin{itemize}
\tightlist
\item
  \textbf{ROC AUC} : \textasciitilde90\% (test)
\item
  \textbf{Recall} : détecte \textasciitilde7 churners sur 10
\item
  \textbf{Precision} : \textasciitilde2 alertes sur 3 justifiées
\end{itemize}
\end{block}

\begin{block}{Variables prédictives}
\phantomsection\label{variables-pruxe9dictives}
\begin{enumerate}
\tightlist
\item
  \textbf{Customer Value}
\item
  \textbf{Complains}
\item
  \textbf{Status}
\item
  \textbf{Subscription Length}
\end{enumerate}
\end{block}
\end{frame}

\begin{frame}{Recommandations métier}
\phantomsection\label{recommandations-muxe9tier}
\begin{block}{Actions prioritaires}
\phantomsection\label{actions-prioritaires}
\begin{enumerate}
\tightlist
\item
  \textbf{Plaintes} → intervention immédiate
\item
  \textbf{Comptes inactifs} → réactivation ciblée
\item
  \textbf{Faible valeur} → offres personnalisées
\end{enumerate}
\end{block}

\begin{block}{KPIs de suivi}
\phantomsection\label{kpis-de-suivi}
\begin{itemize}
\tightlist
\item
  Taux de rétention des clients ciblés
\item
  Coût par client retenu vs CLV
\item
  Drift monitoring du modèle
\end{itemize}
\end{block}
\end{frame}

\begin{frame}{Merci}
\phantomsection\label{merci}
\vspace{1cm}
\centering

\textbf{Questions ?}

\vspace{0.5cm}

\emph{Master MECEN --- Classification Supervisée}

\vspace{1cm}
\footnotesize

Réalisé avec R, tidymodels et Quarto
\end{frame}



\end{document}
